\begin{tip}
Elige el lado del corte que tenga \textbf{menos cargas y menos incógnitas}. No existe premio por escoger siempre el lado izquierdo.
\end{tip}

% ============================================================
\safechapter{Equilibrio interno y leyes de esfuerzos}
% ============================================================

\safesection{Rebanada diferencial}
Consideremos una barra recta sometida a cargas distribuidas:
\[
q_x(x),\qquad q_y(x).
\]
Tomamos una rebanada de longitud $\dd x$.

En una cara actúan
\[
N_x(x),\qquad C_y(x),\qquad M_z(x)
\]
y en la otra
\[
N_x+\dd N_x,\qquad C_y+\dd C_y,\qquad M_z+\dd M_z.
\]

\safesection{Equilibrio axial}
Aplicando equilibrio en $x$:
\[
N_x+\dd N_x-N_x+q_x\dd x=0.
\]
Dividiendo por $\dd x$:
\[
\frac{\dd N_x}{\dd x}=-q_x(x).
\]

\safesection{Equilibrio transversal}
Con la convención de signos del curso:
\[
\frac{\dd C_y}{\dd x}=-q_y(x).
\]

\safesection{Equilibrio de momentos}
Despreciando términos de segundo orden:
\[
\frac{\dd M_z}{\dd x}=C_y(x),
\]
y, combinando:
\[
\frac{\dd^2M_z}{\dd x^2}=-q_y(x).
\]

\safesection{Interpretación geométrica}
De
\[
\frac{\dd M}{\dd x}=V
\]
se deduce:
\flowbox{$V$ = pendiente del diagrama de momentos}

Por tanto, $V=0$ indica un extremo local de $M$.

Además:
\[
\frac{\dd V}{\dd x}=-q.
\]
Así:
\flowbox{$q$ = menos la pendiente del diagrama de cortantes}

\safesection{Regla de grados}
Si $q(x)$ es un polinomio de grado $n$, entonces, dentro de un tramo sin cargas concentradas:
\[
q(x)\sim x^n,
\qquad
V(x)\sim x^{n+1},
\qquad
M(x)\sim x^{n+2}.
\]

\begin{table}[htbp]
\centering
\caption{Relación entre el grado de la carga distribuida y la forma de las leyes de cortante y momento.}
\label{tab:regla-grados-qvm}
\begin{tabular}{ccc}
\toprule
Carga $q$ & Cortante $V$ & Momento $M$\\
\midrule
$0$ & constante & lineal\\
constante & lineal & parabólico\\
lineal & parabólico & cúbico\\
cuadrático & cúbico & grado 4\\
\bottomrule
\end{tabular}
\end{table}

\begin{tip}
Esta tabla permite dibujar muchas leyes \textbf{a estima} sin resolver completamente las ecuaciones.
\end{tip}

\safesection{Discontinuidades}
Una carga puntual transversal produce un salto en $V$. Un momento puntual produce un salto en $M$. Una carga repartida finita no produce saltos: cambia la pendiente.

\[
\text{fuerza puntual}\longrightarrow\text{salto en cortante}
\]
\[
\text{momento puntual}\longrightarrow\text{salto en momento}
\]

\safesection{Condiciones en puntos especiales}

\safesubsection{Extremo libre}
Si no hay carga aplicada justo en el extremo:
\[
N=0,\qquad V=0,\qquad M=0.
\]

\safesubsection{Rótula interna}
Una rótula no transmite momento:
\[
M=0.
\]
